\documentclass{article}
\usepackage[utf8]{inputenc}
\usepackage{amsmath,amssymb}
\usepackage[most]{tcolorbox}
\usepackage{geometry}
\geometry{margin=2.5cm}

\newcommand{\step}[1]{\par\noindent\hangindent=3.5em\hangafter=1%
  \makebox[3.5em][l]{\textbf{#1.}}}
\newcommand{\steptag}[1]{\hfill{\small\textsf{[#1]}}}

\newtcolorbox{proofbox}[1]{
  colback=gray!5, colframe=gray!60,
  fonttitle=\bfseries, title={#1},
  breakable, enhanced,
  left=4pt, right=4pt, top=4pt, bottom=4pt,
  before skip=10pt, after skip=10pt
}

\begin{document}

\begin{proofbox}{Parameterized polar-retraction transport: there exist C\_p...}
\step{1} Parameterized polar-retraction transport: there exist C\_pol\textasciicircum{}0 >= 1 and kappa\_pol\textasciicircum{}0 in (0, 1/2] such that, for every def-stage1-polar-witness-data tuple W with C\_pol >= C\_pol\textasciicircum{}0 and 0 < kappa\_pol <= kappa\_pol\textasciicircum{}0, for every finite-dimensional exact-unit epsilon\_r-C*-algebra and every delta > 0 with C\_pol*(epsilon\_r + delta) <= kappa\_pol, the map Pi\_delta: calU x B\textasciicircum{}\{calH\}\_delta(J) -> calX, Pi\_delta(U, H) = U bold-dot H, is a C\textasciicircum{}1 diffeomorphism onto the open set S\_delta := Pi\_delta(calU x B\textasciicircum{}\{calH\}\_delta(J)), its inverse (u\_delta, h\_delta): S\_delta -> calU x B\textasciicircum{}\{calH\}\_delta(J) satisfies X = u\_delta(X) bold-dot h\_delta(X), u\_delta(U) = U, and h\_delta(U) = J for every X in S\_delta and U in calU, and calU\_\{delta - C\_pol*(epsilon\_r*delta + delta\textasciicircum{}2)\} subseteq S\_delta subseteq calU\_\{delta + C\_pol*(epsilon\_r*delta + delta\textasciicircum{}2)\}. \steptag{validated} \\[6pt]
\step{1.1} By lem-stage1-polar-retraction, fix universal constants C\_* >= 1 and kappa\_* in (0, 1/2] witnessing that lemma. Define C\_pol\textasciicircum{}0 := C\_* and kappa\_pol\textasciicircum{}0 := kappa\_*; these are universal and have the ranges required at the root. \steptag{validated} \\[4pt]
\step{1.2} Fix an arbitrary def-stage1-polar-witness-data tuple W and arbitrary algebra and delta satisfying the root hypotheses. The exact unit J obeys J bold-dot J = J and ||J|| = 1; applying the product-norm axiom of def-epsilon-cstar-algebra to J,J gives 1 = ||J bold-dot J|| <= (1 + epsilon\_r)||J||\textasciicircum{}2 = 1 + epsilon\_r, hence epsilon\_r >= 0. Thus t := epsilon\_r + delta > 0 and A := epsilon\_r*delta + delta\textasciicircum{}2 = delta*t >= 0. Since C\_* <= C\_pol and t >= 0, the assumed guard gives C\_*t <= C\_pol*t <= kappa\_pol <= kappa\_*, which is exactly the guard needed for lem-stage1-polar-retraction at (C\_*, kappa\_*). \steptag{validated} \\[4pt]
\step{1.3} Apply lem-stage1-polar-retraction with its fixed constants (C\_*, kappa\_*) to the arbitrary algebra and delta, using the guard from the preceding step. With S\_delta taken to be the image Pi\_delta(calU x B\textasciicircum{}\{calH\}\_delta(J)), it gives that S\_delta is open, Pi\_delta is a C\textasciicircum{}1 diffeomorphism onto S\_delta, the inverse (u\_delta,h\_delta) has X = u\_delta(X) bold-dot h\_delta(X), u\_delta(U) = U, and h\_delta(U) = J, and it gives the base sandwich calU\_\{delta - C\_*A\} subseteq S\_delta subseteq calU\_\{delta + C\_*A\}. \steptag{validated} \\[4pt]
\step{1.4} Because C\_pol >= C\_* and A >= 0, delta - C\_pol*A <= delta - C\_*A and delta + C\_*A <= delta + C\_pol*A. By def-approximate-unitary-space, if r <= s then calU\_r subseteq calU\_s: the right-inverse condition is unchanged and ||X\textasciicircum{}dagger bold-dot X - J|| < 2r implies the same strict inequality with 2s. Hence calU\_\{delta - C\_pol*A\} subseteq calU\_\{delta - C\_*A\} subseteq S\_delta subseteq calU\_\{delta + C\_*A\} subseteq calU\_\{delta + C\_pol*A\}. Together with the unchanged open-set, diffeomorphism, and inverse conclusions from lem-stage1-polar-retraction, this proves the root statement for the arbitrary W, algebra, and delta. \steptag{validated} \\[4pt]
\end{proofbox}

\end{document}
